\chapter{월 사용료}

\begin{summarybox}{한눈에}
본사에 매월 지불하는 \textbf{기관 월 사용료}를 청구·결제·이력 조회합니다. 매월 1일 새벽 자동 발행, 그 달 말일까지 결제, 다음 달 1일 미결제 시 기관 관리 기능 자동 정지.
\end{summarybox}

\kfShot{13-billing}{월 사용료 — 청구 이력 + 산정 미리보기 + 결제}

\section*{산정 방식}

\begin{itemize}[leftmargin=1.2em]
  \item \textbf{0\,$\sim$\,20명}: 기본료 200,000원 고정 (본사가 기관별 감면 가능)
  \item \textbf{21\,$\sim$\,50명}: 기본료 + (가장 등록일수 짧은 (학생수-20)명의 등록일수 합 × 500원)
  \item 등록일수 = 그 달 안에 status='active' 였던 일수 (가입 중간이면 가입일\,$\sim$\,말일)
\end{itemize}

예시) 24명 등록 시: 기본료 200,000 + (등록일수 짧은 4명의 일수 합 × 500원)

\section*{화면 구조}

\begin{screenbox}{청구 이력 + 실시간 산정 미리보기}
\textbf{이번 달 예상 사용료 (실시간)}
\noindent
\begin{tabular}{@{}p{4cm}p{4cm}@{}}
활성 학생 수 & N 명 \\
기본료 & 200,000원 \\
초과 학생 & N 명 \\
초과 등록일수 합 & N 일 × 500원 = N원 \\
\midrule
\textbf{합계} & \textbf{N원} \\
\bottomrule
\end{tabular}

\medskip
\textbf{청구 이력}
\noindent
\begin{tabular}{@{}p{1.8cm}p{2cm}p{2.5cm}p{2cm}p{2cm}p{2cm}@{}}
\toprule
\textbf{월} & \textbf{기본료} & \textbf{초과/일수} & \textbf{합계} & \textbf{상태} & \textbf{결제} \\
\midrule
2026-05 & 200,000 & 5명 / 80일 & 240,000 & 결제 대기 & \activetab{결제} \\
2026-04 & 200,000 & 0 / 0 & 200,000 & 결제 완료 & — \\
\bottomrule
\end{tabular}
\end{screenbox}

\section*{여기서 할 수 있는 일}
\begin{itemize}[leftmargin=1.2em]
  \item 이번 달 \textbf{실시간 예상 사용료} 미리 계산 (학생 수 변동에 따라 갱신)
  \item 발행된 청구서 \textbf{결제} (mock — 정식 토스 연동은 추후)
  \item \textbf{청구 이력} 조회 (월별 기본료·초과·합계·상태)
  \item 정지 상태 시 결제 → 즉시 해제
\end{itemize}

\section*{자주 쓰는 흐름}
\begin{enumerate}
  \item 매월 1일 → 청구서 자동 발행 + 화면 상단에 빨간/노란 배너
  \item 결제 마감 7일 전부터 노란 ``결제 임박'' 배너
  \item ``결제'' 버튼 → mock 결제 → 즉시 paid 처리
  \item 마감 지난 후 새벽 3시에 자동 정지 → 결제하면 즉시 해제
\end{enumerate}

\begin{tipbox}{알아두면 좋은 점}
\begin{itemize}[leftmargin=1.2em]
  \item 정지 상태에서도 \textbf{이 메뉴(``월 사용료'') 와 자몽 충전 / 기관 설정 / 로그인}은 가능
  \item 학생 관리·콘텐츠 등 다른 기관 운영은 차단됨
  \item 본사가 기관별로 \textbf{기본료 감면} 가능 — 협의 후 본사가 설정
  \item 학생 학습은 미결제와 무관하게 정상 진행 (학생이 학원 정책으로 피해보지 않도록)
\end{itemize}
\end{tipbox}
